\subsection{Real Spectral Data via Complexification}

\begin{lemma} \label{lem:complexification-conjugate-generalized-eigenspaces} \hlt{Conjugation Transfers Generalized Eigenspaces}\\
Let $\mathcal{V}$ be a real vector space and let $T\in\LL(\mathcal{V})$. For $\mu\in\mathbb{C}$, set $S_{\mu}\defeq T_{\mathbb{C}}-\mu I_{\mathcal{V}_{\mathbb{C}}}$.\\
Then $(\forall\lambda\in\mathbb{C})(\forall k\in\mathbb{Z}_{\geq0})(\forall\mathbf{u},\mathbf{v}\in\mathcal{V})\,\bigl(S_{\lambda}^k(\mathbf{u}+i\mathbf{v})=\mathbf{0}\iff S_{\overline{\lambda}}^k(\mathbf{u}-i\mathbf{v})=\mathbf{0}\bigr)$.
\begin{proof}
Define $\Gamma:\mathcal{V}_{\mathbb{C}}\to\mathcal{V}_{\mathbb{C}}$ by $\Gamma(\mathbf{a}+i\mathbf{b})\defeq\mathbf{a}-i\mathbf{b}$. By Definition \ref{def:complexification-real-vector-space}, $(\forall\mathbf{x},\mathbf{y}\in\mathcal{V}_{\mathbb{C}})\,\Gamma(\mathbf{x}+\mathbf{y})=\Gamma\mathbf{x}+\Gamma\mathbf{y}$.\\
The same definition implies $(\forall\alpha\in\mathbb{C})(\forall\mathbf{x}\in\mathcal{V}_{\mathbb{C}})\,\Gamma(\alpha\mathbf{x})=\overline{\alpha}\Gamma\mathbf{x}$ and $\Gamma^2=I_{\mathcal{V}_{\mathbb{C}}}$, so $\Gamma$ is a conjugate-linear bijection.\\
By Definition \ref{def:complexification-real-operator}, $\Gamma T_{\mathbb{C}}=T_{\mathbb{C}}\Gamma$. Thus $(\forall\mu\in\mathbb{C})\,\Gamma S_{\mu}=S_{\overline{\mu}}\Gamma$. The corresponding power identity holds at exponent $0$.\\
If it holds at $r\in\mathbb{Z}_{\geq0}$, then it holds at $r+1$ by composition with $\Gamma S_{\lambda}=S_{\overline{\lambda}}\Gamma$; $(\forall r\in\mathbb{Z}_{\geq0})(\forall\mathbf{x}\in\mathcal{V}_{\mathbb{C}})\,\Gamma(S_{\lambda}^r\mathbf{x})=S_{\overline{\lambda}}^r\Gamma\mathbf{x}$ by induction.
\begin{enumerate}[label=\roman*.]
\setlength{\itemsep}{0pt}
\item[$(\Rightarrow)$] If $S_{\lambda}^k(\mathbf{u}+i\mathbf{v})=\mathbf{0}$, then $S_{\overline{\lambda}}^k(\mathbf{u}-i\mathbf{v})=\Gamma\mathbf{0}=\mathbf{0}$ by the power identity.
\item[$(\Leftarrow)$] If $S_{\overline{\lambda}}^k(\mathbf{u}-i\mathbf{v})=\mathbf{0}$, then $S_{\lambda}^k(\mathbf{u}+i\mathbf{v})=\Gamma\mathbf{0}=\mathbf{0}$ by the power identity with $\overline{\lambda}$ in place of $\lambda$.
\end{enumerate}
\hlt{Structural Commentary}: Conjugate-linearity intertwines every shifted power with its conjugate shift.\\
\hlt{Key Insight}: The involution $\Gamma$ transports each shifted kernel to its conjugate kernel.
\end{proof}
\end{lemma}

\begin{corollary} \label{cor:complexification-conjugate-eigenvalue-multiplicities} \hlt{Conjugate Eigenvalues Have Equal Multiplicity}\\
Let $\mathcal{V}\neq\{\mathbf{0}\}$ be finite-dimensional over $\mathbb{R}$ and let $T\in\LL(\mathcal{V})$. For $\lambda\in\operatorname{spec}(T_{\mathbb{C}})$:
\begin{enumerate}[label=(\roman*)]
\setlength{\itemsep}{0pt}
\item $\overline{\lambda}\in\operatorname{spec}(T_{\mathbb{C}})$;
\item $\dim_{\mathbb{C}}\mathcal{G}_{\lambda}(T_{\mathbb{C}})=\dim_{\mathbb{C}}\mathcal{G}_{\overline{\lambda}}(T_{\mathbb{C}})$;
\item $\lambda$ and $\overline{\lambda}$ have the same algebraic multiplicity.
\end{enumerate}
\begin{proof}
For $\mu\in\mathbb{C}$, set $S_{\mu}\defeq T_{\mathbb{C}}-\mu I_{\mathcal{V}_{\mathbb{C}}}$. Define $\Gamma:\mathcal{V}_{\mathbb{C}}\to\mathcal{V}_{\mathbb{C}}$ by $\Gamma(\mathbf{u}+i\mathbf{v})\defeq\mathbf{u}-i\mathbf{v}$. By Definition \ref{def:complexification-real-vector-space}, $(\forall\mathbf{x},\mathbf{y}\in\mathcal{V}_{\mathbb{C}})\,\Gamma(\mathbf{x}+\mathbf{y})=\Gamma\mathbf{x}+\Gamma\mathbf{y}$.\\
The same definition implies $(\forall\alpha\in\mathbb{C})(\forall\mathbf{x}\in\mathcal{V}_{\mathbb{C}})\,\Gamma(\alpha\mathbf{x})=\overline{\alpha}\Gamma\mathbf{x}$ and $\Gamma^2=I_{\mathcal{V}_{\mathbb{C}}}$, so $\Gamma$ is a conjugate-linear bijection.
\begin{enumerate}[label=(\roman*)]
\setlength{\itemsep}{0pt}
\item Choose a nonzero eigenvector $\mathbf{z}=\mathbf{u}+i\mathbf{v}$ for $\lambda$. By Lemma \ref{lem:complexification-conjugate-generalized-eigenspaces} with $k=1$, $S_{\overline{\lambda}}\Gamma\mathbf{z}=\mathbf{0}$.\\
The vector $\Gamma\mathbf{z}$ is nonzero by $\Gamma^2=I_{\mathcal{V}_{\mathbb{C}}}$, so $\overline{\lambda}\in\operatorname{spec}(T_{\mathbb{C}})$ by Definition \ref{def:eigenvalue-eigenvector}.
\item Set $n\defeq\dim_{\mathbb{C}}\mathcal{V}_{\mathbb{C}}$. By Theorem \ref{thm:generalized-eigenspace-stabilized-kernel}(i), the generalized eigenspaces are $\ker(S_{\lambda}^n)$ and $\ker(S_{\overline{\lambda}}^n)$.\\
By Lemma \ref{lem:complexification-conjugate-generalized-eigenspaces} with $k=n$ and $\Gamma^2=I_{\mathcal{V}_{\mathbb{C}}}$, $\Gamma$ restricts to a conjugate-linear bijection $\mathcal{G}_{\lambda}(T_{\mathbb{C}})\to\mathcal{G}_{\overline{\lambda}}(T_{\mathbb{C}})$.\\
The space $\mathcal{G}_{\lambda}(T_{\mathbb{C}})$ is finite-dimensional by Theorem \ref{thm:generalized-eigenspace-stabilized-kernel}(ii) and Corollary \ref{cor:finite-dimensional-spaces-and-subspaces-have-bases}(ii).\\
Set $d\defeq\dim_{\mathbb{C}}\mathcal{G}_{\lambda}(T_{\mathbb{C}})$ and choose a basis $(\mathbf{z}_1,\ldots,\mathbf{z}_d)$ by Corollary \ref{cor:finite-dimensional-spaces-and-subspaces-have-bases}(i).\\
If $\sum\limits_{j=1}^{d}a_j\Gamma\mathbf{z}_j=\mathbf{0}$, then $\sum\limits_{j=1}^{d}\overline{a_j}\mathbf{z}_j=\mathbf{0}$ by conjugate-linearity and $\Gamma^2=I_{\mathcal{V}_{\mathbb{C}}}$. Hence $(\forall j\in\{1,\ldots,d\})\,a_j=0$ by independence.\\
Fix $\mathbf{y}\in\mathcal{G}_{\overline{\lambda}}(T_{\mathbb{C}})$ and choose $b_1,\ldots,b_d\in\mathbb{C}$ satisfying $\Gamma\mathbf{y}=\sum\limits_{j=1}^{d}b_j\mathbf{z}_j$ by spanning.\\
Then $\mathbf{y}=\sum\limits_{j=1}^{d}\overline{b_j}\Gamma\mathbf{z}_j$ by conjugate-linearity, $\Gamma^2=I_{\mathcal{V}_{\mathbb{C}}}$. Thus $(\Gamma\mathbf{z}_1,\ldots,\Gamma\mathbf{z}_d)$ is basis of $\mathcal{G}_{\overline{\lambda}}(T_{\mathbb{C}})$, its dimension is $d=\dim_{\mathbb{C}}\mathcal{G}_{\lambda}(T_{\mathbb{C}})$.
\item The algebraic multiplicities are the dimensions in (ii) by Proposition \ref{prop:eigenvalue-algebraic-multiplicity-generalized-eigenspace}.
\end{enumerate}
\hlt{Structural Commentary}: Kernel transport constructs an image basis; the multiplicity formula transfers its dimension equality.\\
\hlt{Key Insight}: Conjugation preserves generalized-eigenspace dimension.
\end{proof}
\end{corollary}

\begin{corollary} \label{cor:odd-dimensional-real-operator-has-eigenvalue} \hlt{Odd-Dimensional Real Operators Have Eigenvalues}\\
Let $\mathcal{V}\neq\{\mathbf{0}\}$ be a finite-dimensional real vector space with $\dim_{\mathbb{R}}\mathcal{V}$ odd. Then $(\forall T\in\LL(\mathcal{V}))\,\operatorname{spec}(T)\neq\varnothing$.
\begin{proof}
Fix $T\in\LL(\mathcal{V})$ and set $n\defeq\dim_{\mathbb{R}}\mathcal{V}$. Then $\dim_{\mathbb{C}}\mathcal{V}_{\mathbb{C}}=n$ by Proposition \ref{prop:complexification-vector-space-basis}(iii).\\
The algebraic multiplicities of the eigenvalues of $T_{\mathbb{C}}$ sum to $n$ by Proposition \ref{prop:sum-of-complex-eigenvalue-multiplicities}.\\
If $\operatorname{spec}(T_{\mathbb{C}})\cap\mathbb{R}=\varnothing$, then the spectrum partitions into pairs $\{\lambda,\overline{\lambda}\}$.\\
Each pair has total equal to twice its common algebraic multiplicity by Corollary \ref{cor:complexification-conjugate-eigenvalue-multiplicities}(i),(iii). Thus $n$ would be even, a contradiction.\\
Thus $T_{\mathbb{C}}$ has a real eigenvalue, which is an eigenvalue of $T$ by Proposition \ref{prop:real-eigenvalues-complexification}.\\
\hlt{Structural Commentary}: Parity compares the total complex multiplicity with the even totals of nonreal conjugate pairs.\\
\hlt{Key Insight}: An odd multiplicity sum cannot consist only of conjugate pairs.
\end{proof}
\end{corollary}

\begin{theorem} \label{thm:real-operator-characteristic-polynomial-properties} \hlt{Characteristic Polynomial over $\mathbb{R}$}\\
Let $\mathcal{V}\neq\{\mathbf{0}\}$ be a finite-dimensional real vector space and let $T\in\LL(\mathcal{V})$. Let $q_T$ be the characteristic polynomial of $T$.
\begin{enumerate}[label=(\roman*)]
\setlength{\itemsep}{0pt}
\item $q_T\in\mathbb{R}[z]$.
\item $\deg q_T=\dim_{\mathbb{R}}\mathcal{V}$.
\item The real zeros of $q_T$ are precisely the eigenvalues of $T$.
\item $q_T(T)=0$.
\end{enumerate}
\begin{proof}
\begin{enumerate}[label=(\roman*)]
\setlength{\itemsep}{0pt}
\item The relation $q_T\in\mathbb{R}[z]$ holds by Definition \ref{def:operator-characteristic-polynomial}.
\item By Proposition \ref{prop:operator-characteristic-polynomial-degree-roots}(i), $\deg q_T=\dim_{\mathbb{R}}\mathcal{V}$.
\item By Proposition \ref{prop:operator-characteristic-polynomial-degree-roots}(ii), the zeros of $q_T$ in $\mathbb{R}$ are precisely the eigenvalues of $T$.
\item Choose a basis $\beta$ of $\mathcal{V}$ by Corollary \ref{cor:finite-dimensional-spaces-and-subspaces-have-bases}(i), and set $\beta_{\mathbb{C}}\defeq\operatorname{cpx}(\beta)$. The list $\beta_{\mathbb{C}}$ is a complex basis by Proposition \ref{prop:complexification-vector-space-basis}(ii).\\
The identity $[T_{\mathbb{C}}]_{\beta_{\mathbb{C}}\leftarrow\beta_{\mathbb{C}}}=[T]_{\beta\leftarrow\beta}$ holds by Proposition \ref{prop:complexification-operator-basic-properties}(iii). Hence $q_{T_{\mathbb{C}}}=q_T$ by Definition \ref{def:operator-characteristic-polynomial}.\\
By Theorem \ref{thm:cayley-hamilton-operator}, $q_T(T_{\mathbb{C}})=q_{T_{\mathbb{C}}}(T_{\mathbb{C}})=0$. Fix $\mathbf{u}\in\mathcal{V}$.\\
By (i), Proposition \ref{prop:complexification-operator-basic-properties}(ii), and linearity of polynomial evaluation, $\mathbf{0}=q_T(T_{\mathbb{C}})(\mathbf{u}+i\mathbf{0})=q_T(T)\mathbf{u}+i\mathbf{0}$.\\
Thus $q_T(T)\mathbf{u}=\mathbf{0}$ by Definition \ref{def:complexification-real-vector-space}(iv), so $q_T(T)=0$ by the arbitrariness of $\mathbf{u}$.
\end{enumerate}
\hlt{Structural Commentary}: Definition and root-degree properties settle (i)--(iii); unchanged coordinates transfer Cayley--Hamilton for (iv).\\
\hlt{Key Insight}: Complexification preserves the real coordinate matrix of $T$.
\end{proof}
\end{theorem}

The identity $q_{T_{\mathbb{C}}}=q_T$ holds by Proposition \ref{prop:complexification-operator-basic-properties}(iii) and Definition \ref{def:operator-characteristic-polynomial}. If $\lambda\in\operatorname{spec}(T_{\mathbb{C}})\setminus\mathbb{R}$ has algebraic multiplicity $m$, then $\overline{\lambda}$ has multiplicity $m$ by Corollary \ref{cor:complexification-conjugate-eigenvalue-multiplicities}(iii). Thus the paired factor is $(z-\lambda)^m(z-\overline{\lambda})^m=\bigl(z^2-2\operatorname{Re}(\lambda)z+\abs{\lambda}^2\bigr)^m$.\\
Figure \ref{fig:real-characteristic-conjugate-pairs} records this real quadratic contribution.

\begin{illustrationfigure}[H]
  \centering
  \begin{tikzpicture}[
  >=Stealth,
  every node/.style={font=\small},
  box/.style={draw=gray!55, rounded corners=3pt, inner xsep=6.5pt, inner ysep=4.5pt, align=center},
  arrow/.style={->, line width=0.85pt, blue!70!black}
]
  \node[box] (complex) at (0,0) {$(z-\lambda)^m(z-\overline{\lambda})^m$};
  \node[box] (real) at (3.87,0) {$\bigl(z^2-2\operatorname{Re}(\lambda)z+\abs{\lambda}^2\bigr)^m$};
  \node[box] (poly) at (7.17,0) {$q_T\in\mathbb{R}[z]$};

  \draw[arrow] (complex) -- (real);
  \draw[arrow] (real) -- (poly);
\end{tikzpicture}

  \caption{Conjugate factors pair into real quadratics.}
  \label{fig:real-characteristic-conjugate-pairs}
\end{illustrationfigure}

\begin{corollary} \label{cor:real-characteristic-polynomial-multiple-minimal-polynomial} \hlt{$q_T$ Is a Multiple of $\mu_T$ over $\mathbb{R}$}\\
Let $\mathcal{V}\neq\{\mathbf{0}\}$ be a finite-dimensional real vector space and let $T\in\LL(\mathcal{V})$. Then:
\begin{enumerate}[label=(\roman*)]
\setlength{\itemsep}{0pt}
\item $\mu_T$ divides $q_T$ in $\mathbb{R}[z]$;
\item $\deg\mu_T\leq\dim_{\mathbb{R}}\mathcal{V}$.
\end{enumerate}
\begin{proof}
By Theorem \ref{thm:real-operator-characteristic-polynomial-properties}(iv), $q_T(T)=0$.
\begin{enumerate}[label=(\roman*)]
\setlength{\itemsep}{0pt}
\item The divisibility $\mu_T\mid q_T$ in $\mathbb{R}[z]$ holds by Theorem \ref{thm:operator-annihilator-iff-multiple-minimal-polynomial}.
\item By (i) and Theorem \ref{thm:real-operator-characteristic-polynomial-properties}(ii), $\deg\mu_T\leq\deg q_T=\dim_{\mathbb{R}}\mathcal{V}$.
\end{enumerate}
\hlt{Structural Commentary}: Real Cayley--Hamilton supplies an annihilator; minimal-polynomial divisibility then controls its degree.\\
\hlt{Key Insight}: The least monic annihilator divides $q_T$.
\end{proof}
\end{corollary}

\begin{proposition} \label{prop:minimal-polynomial-complexification} \hlt{Minimal Polynomial Survives Complexification}\\
Let $\mathcal{V}\neq\{\mathbf{0}\}$ be finite-dimensional over $\mathbb{R}$ and let $T\in\LL(\mathcal{V})$. Under $\mathbb{R}[z]\subseteq\mathbb{C}[z]$, $\mu_{T_{\mathbb{C}}}=\mu_T$.
\begin{proof}
\begin{enumerate}[label=(\roman*)]
\setlength{\itemsep}{0pt}
\item \hlt{Extended Annihilation}: Fix $\mathbf{u},\mathbf{v}\in\mathcal{V}$.\\
Then $\mu_T(T_{\mathbb{C}})(\mathbf{u}+i\mathbf{v})=\mathbf{0}$ by $\mu_T(T)=0$, Proposition \ref{prop:complexification-operator-basic-properties}(ii), and linearity of polynomial evaluation.\\
Thus $\mu_T(T_{\mathbb{C}})=0$ by the arbitrariness of $\mathbf{u},\mathbf{v}$, and $\mu_{T_{\mathbb{C}}}$ divides $\mu_T$ in $\mathbb{C}[z]$ by Theorem \ref{thm:operator-annihilator-iff-multiple-minimal-polynomial}.
\item \hlt{Real-Part Descent}: Set $q\defeq\mu_{T_{\mathbb{C}}}=r+is$, where $r,s\in\mathbb{R}[z]$. Since $q$ is monic, $r$ is monic and $\deg r=\deg q$.\\
Fix $\mathbf{u}\in\mathcal{V}$. By Definition \ref{def:operator-minimal-polynomial}, Proposition \ref{prop:complexification-operator-basic-properties}(ii), and linearity of polynomial evaluation, $\mathbf{0}=q(T_{\mathbb{C}})(\mathbf{u}+i\mathbf{0})=r(T)\mathbf{u}+is(T)\mathbf{u}$.\\
Hence $r(T)\mathbf{u}=s(T)\mathbf{u}=\mathbf{0}$ by Definition \ref{def:complexification-real-vector-space}(iv), so $r(T)=s(T)=0$ by the arbitrariness of $\mathbf{u}$.\\
Therefore $\mu_T$ divides $r$ in $\mathbb{R}[z]$ by Theorem \ref{thm:operator-annihilator-iff-multiple-minimal-polynomial}.
\item \hlt{Degree Comparison}: By the first two stages and $\deg r=\deg q$, $\deg q\leq\deg\mu_T\leq\deg r=\deg q$; hence $\deg q=\deg\mu_T$.
\item \hlt{Monic Identification}: Choose $c\in\mathbb{C}^{\times}$ satisfying $\mu_T=cq$ by divisibility and degree equality.\\
Both polynomials are monic, so $c=1$ and $\mu_{T_{\mathbb{C}}}=q=\mu_T$.
\end{enumerate}
\hlt{Structural Commentary}: Extended annihilation and real-part descent force equal degrees before monicity identifies the divisors.\\
\hlt{Key Insight}: The real part of the complex minimal polynomial remains monic and annihilates $T$.
\end{proof}
\end{proposition}
